\chapter{验证AI输出正误的方法}

上一章讨论了 AI 在机器学习项目中经常犯哪些错误。本章进一步回答一个更实际的问题：当 AI 已经生成代码、模型和分析结果之后，我们应当怎样判断它们是否正确？

这里的“验证”不是在流程末尾看一眼测试集分数，也不是让 AI 再解释一次自己的答案。它是一组可以重复执行的检查方法。每种方法从不同角度提出问题：复杂模型是否真的优于简单方法？结果是否符合现实常识？模型遇到边界情况会怎样？错误是否集中在某些时间、区域或类别？只有把这些问题组合起来，评价结果才具有可信度。

本章重点介绍五种实用方法：基准测试、量纲与数量级检验、极端情景测试、反向追问 AI 和可视化验证。它们不要求使用复杂理论，却能够发现许多仅凭模型指标无法发现的问题。

% ----------------------------------------------------------------
\section{基准测试：高级模型是否真的更好}
% ----------------------------------------------------------------

\subsection{为什么必须建立基准}

假设 AI 建立了一个交通流预测模型，并报告测试集上的平均绝对误差为 18。这个数字本身几乎没有意义：18 是否足够小，取决于流量的规模、业务要求，以及更简单的方法能够达到什么水平。

基准模型回答的是一个朴素但关键的问题：

\begin{quote}
如果不用这个复杂模型，只采用最简单、最自然的方法，结果会有多好？
\end{quote}

只有当复杂模型稳定优于合理基准时，增加的建模复杂度才可能具有价值。如果一个深度学习模型还不如“使用上一时刻的观测值”，继续调参往往不是最优先的工作。此时更应该检查任务定义、数据质量、特征构造和评价方式。

\subsection{如何选择有意义的基准}

基准不需要复杂，但必须与任务相匹配。常见选择如下：

\begin{table}[htbp]
\centering
\small
\begin{tabular}{p{0.22\textwidth}p{0.30\textwidth}p{0.38\textwidth}}
\toprule
任务类型 & 可用基准 & 基准回答的问题 \\
\midrule
时间序列预测 & 上一时刻值、历史同时刻均值、移动平均 & 模型是否优于时间序列本身的惯性和周期性 \\
分类任务 & 多数类别、简单规则、逻辑回归 & 复杂模型是否真正识别了额外规律 \\
回归任务 & 全局均值、分组均值、线性回归 & 模型是否优于简单统计关系 \\
排序与推荐 & 热门项目、最近项目、人工规则 & 个性化模型是否优于通用选择 \\
\bottomrule
\end{tabular}
\end{table}

一个项目通常不应只有一个基准。例如，短时交通流预测可以同时比较“上一时刻流量”和“上周同一时刻流量”。前者代表短期惯性，后者代表周期性。模型只有同时优于二者，才能说明它学到了更有用的信息。

\subsection{公平地比较模型与基准}

基准和候选模型必须在相同条件下比较：

\begin{itemize}
  \item 使用相同的训练集、验证集和测试集划分；
  \item 使用相同的评价指标和样本范围；
  \item 不允许候选模型使用预测时不可获得的信息；
  \item 分别报告总体表现和关键场景表现；
  \item 不只比较一次运行中的最好结果。
\end{itemize}

如果复杂模型只在总体平均指标上略好，但在高峰期、少数类别或关键区域明显更差，就不能简单地宣布模型胜出。基准测试的目的不是为复杂模型寻找胜利条件，而是判断复杂模型究竟增加了什么价值。

\subsection{一个基准测试案例}

AI 为某城市建立了一个次日客流预测模型，并声称模型的 RMSE 很低。加入三个基准后，结果如下：

\begin{itemize}
  \item 使用全局平均值：误差较大；
  \item 使用上周同一天客流：已经达到较好结果；
  \item AI 生成的复杂模型：总体误差比“上周同一天”低 4\%，但节假日误差更高。
\end{itemize}

这个结果并不意味着复杂模型毫无价值，但它改变了评价重点。接下来应研究模型为何在节假日失效，以及 4\% 的改善是否值得增加部署和维护成本，而不是仅凭“复杂模型分数最高”结束评价。

% ----------------------------------------------------------------
\section{量纲与数量级检验：结果在现实中是否合理}
% ----------------------------------------------------------------

\subsection{先检查单位，再解释结果}

AI 可以生成语法正确、计算完整的代码，却不一定理解数据中每个数值的真实含义。速度可能以 km/h 或 m/s 表示，距离可能以米或千米表示，时间可能以秒或分钟表示。单位混淆常常不会导致程序报错，却会使结果相差数倍甚至数千倍。

因此，对关键输入和输出都应记录：

\begin{itemize}
  \item 数值代表什么对象；
  \item 使用什么单位；
  \item 合理范围大致是多少；
  \item 数据中的特殊值代表什么；
  \item 转换单位后是否仍与原定义一致。
\end{itemize}

例如，模型预测一段道路的平均速度为 240 km/h，或者预测某站点每分钟出现 50 万名乘客，即使评价代码正常运行，也应立即停止解释模型并检查数据与计算过程。

\subsection{用数量级发现隐藏错误}

数量级检验不要求提前知道精确答案，而是判断结果是否处于合理范围。常见问题包括：

\begin{itemize}
  \item 总量与单条记录的量级混淆；
  \item 百分数和比例混淆，例如将 0.18 当作 18；
  \item 时间窗口错误，例如把每分钟流量当作每小时流量；
  \item 重复连接导致总量成倍增加；
  \item 缺失值填充后使均值或方差异常变化。
\end{itemize}

例如，数据连接前有 10 万条出行记录，连接后变成 900 万条。连接操作可能在技术上完全成功，但数量级已经提示连接键并不唯一。此时应检查每个键的匹配数量，而不是继续训练模型。

\subsection{检查边界、方向和关系}

除了绝对大小，还应检查变量之间应当满足的基本关系：

\begin{description}
  \item[边界检查] 概率应在 0 到 1 之间，持续时间不应为负，类别编码不应出现未知取值；
  \item[方向检查] 在其他条件相近时，极端拥堵通常不应对应异常高速度；
  \item[组成关系检查] 各类别数量之和应与总数量一致；
  \item[时间关系检查] 用于预测的特征必须早于预测时点；
  \item[空间关系检查] 地理坐标和距离应落在研究区域及合理范围内。
\end{description}

这些规则可以直接写成断言或数据质量检查。AI 可以协助生成检查代码，但哪些关系必须成立，需要由理解任务的人决定。

\subsection{建立“合理性检查表”}

每个项目都可以维护一张简短的合理性检查表。例如：

\begin{table}[htbp]
\centering
\small
\begin{tabular}{p{0.25\textwidth}p{0.30\textwidth}p{0.35\textwidth}}
\toprule
检查对象 & 预期范围或关系 & 发现异常后的动作 \\
\midrule
预测速度 & 0--道路限速附近 & 检查单位、异常值与输出变换 \\
训练与测试时间 & 训练数据早于测试数据 & 停止训练并重新划分数据 \\
连接后的记录数 & 与连接关系预期一致 & 检查连接键唯一性 \\
分类概率 & 0--1，且与类别定义一致 & 检查模型输出与后处理代码 \\
\bottomrule
\end{tabular}
\end{table}

合理性检查表的价值在于把领域知识转化为可以执行的评价条件，而不是等看到奇怪结果后才临时判断。

% ----------------------------------------------------------------
\section{极端情景测试：模型在边界条件下会怎样}
% ----------------------------------------------------------------

\subsection{平均表现为什么不够}

测试集通常以常见样本为主。因此，一个模型即使总体指标很好，也可能在少见但重要的场景中严重失效。例如，交通预测模型可能在普通工作日表现稳定，却在暴雨、节假日或大型活动期间产生明显错误。

极端情景测试也称压力测试。它不只问“模型平均有多准”，还问：

\begin{quote}
当输入偏离常规情况，或者项目最担心的事件发生时，模型会给出什么结果？
\end{quote}

\subsection{从任务中列出关键情景}

极端情景不一定是数学意义上的最大值或最小值，而是对真实使用具有重要影响的边界条件。可以从以下方向设计：

\begin{itemize}
  \item \textbf{时间变化：}高峰、深夜、节假日、季节转换；
  \item \textbf{环境变化：}恶劣天气、突发事件、大型活动；
  \item \textbf{数据异常：}缺失字段、传感器中断、异常值、延迟数据；
  \item \textbf{对象变化：}新用户、新区域、新设备、低频类别；
  \item \textbf{规模变化：}输入量突然增大、并发请求增加、处理时间变长。
\end{itemize}

选择情景时，应优先考虑“发生概率不一定最高，但一旦出错后果明显”的情况。

\subsection{先写预期，再运行测试}

极端情景测试不一定要求提前知道精确预测值，但应提前写出合理方向和允许范围。例如：

\begin{itemize}
  \item 当一半传感器缺失时，系统应发出警告，而不是继续输出高置信度预测；
  \item 当输入超出训练范围时，模型应被标记为低可信，而不是自动外推；
  \item 当某类别完全没有训练样本时，系统不应声称能够可靠识别；
  \item 当节假日规律与工作日不同，应单独报告节假日误差。
\end{itemize}

先写预期能够避免在看到结果后为模型寻找解释。它也使极端测试更接近实验，而不是随意观察。

\subsection{记录失败方式，而不只记录分数}

极端测试应记录模型如何失败：

\begin{itemize}
  \item 性能逐渐下降，还是突然崩溃；
  \item 输出是否超过合理范围；
  \item 置信度是否随着错误增加而降低；
  \item 系统是否能够识别异常输入；
  \item 是否存在可用的基准方案或人工处理方式。
\end{itemize}

评价的目标不是证明模型能应对所有情况，而是明确它在哪些条件下仍可使用，在哪些条件下必须停止依赖。

% ----------------------------------------------------------------
\section{反向追问AI：让AI帮助发现自己的问题}
% ----------------------------------------------------------------

\subsection{从“生成答案”切换到“寻找错误”}

AI 在第一次回答时，通常倾向于完成任务并给出连贯方案。验证阶段需要改变它的任务：不再要求继续完善答案，而是要求主动寻找可能使答案错误的原因。

例如，可以使用以下 Prompt：

\begin{quote}
请不要继续优化这个模型。假设当前结论可能是错误的，列出至少五个最可能的问题，并按发生可能性排序。对每个问题说明应检查什么证据。
\end{quote}

也可以要求 AI 扮演更明确的审查角色：

\begin{quote}
请以严格的机器学习论文审稿人身份审查这项实验。重点检查数据泄露、基准选择、评价指标、分组表现和结论是否超过证据。
\end{quote}

\subsection{让AI提出可执行的验证动作}

只让 AI “指出问题”容易得到泛泛的风险列表。更有效的要求是让它把每个怀疑转化为具体检查：

\begin{itemize}
  \item 需要查看哪个字段、样本或中间结果；
  \item 应构造什么测试数据；
  \item 预期看到什么现象；
  \item 什么结果会推翻当前结论；
  \item 修正后需要重新运行哪些步骤。
\end{itemize}

例如，AI 怀疑存在时间泄露时，应进一步要求它列出可能泄露的特征，并为每个特征说明生成时间和预测时点，而不是只接受“注意避免数据泄露”这一句提醒。

\subsection{使用独立上下文减少共同错误}

如果可能，可以在新的对话中让 AI 审查结果，只提供任务定义、代码和输出，而不提供原对话中的解释与辩护。这样可以减少审查过程沿用第一次回答的假设。

对于重要任务，还可以让不同模型、不同工具或人工审查者分别检查，再比较它们发现的问题。不同检查者的一致意见值得重视，不一致之处则提示需要进一步核验。

\subsection{AI审查仍然不是独立证据}

AI 可以扩大检查范围，却不能代替实际验证。生成者和审查者可能共享相同的错误知识，也可能都没有注意到数据定义本身有误。因此，反向追问应当用于提出检查方向，而不是直接证明结果正确。

% ----------------------------------------------------------------
\section{可视化验证：让错误看得见}
% ----------------------------------------------------------------

\subsection{为什么图形能够发现指标遗漏的问题}

单个评价指标会压缩大量信息。两个模型可能具有相同的平均误差，却以完全不同的方式犯错：一个在所有样本上略有偏差，另一个在少数场景中出现巨大错误。可视化能够展示误差发生的位置、时间和模式。

\subsection{根据任务选择验证图形}

\begin{table}[htbp]
\centering
\small
\begin{tabular}{p{0.22\textwidth}p{0.30\textwidth}p{0.38\textwidth}}
\toprule
任务 & 推荐图形 & 重点观察 \\
\midrule
时间序列预测 & 真实值与预测值曲线、误差随时间变化图 & 高峰、转折点和异常日期是否持续失效 \\
回归任务 & 真实值与预测值散点图、残差图 & 是否系统性高估或低估，误差是否随数值增大 \\
分类任务 & 混淆矩阵、概率分布、错误案例表 & 哪些类别容易混淆，置信度是否合理 \\
空间任务 & 误差地图、分区表现图 & 错误是否集中在特定区域 \\
数据处理 & 处理前后分布图、缺失率图 & 清洗是否改变了原有结构 \\
\bottomrule
\end{tabular}
\end{table}

\subsection{图形必须服务于具体问题}

验证图形不是为了让报告更丰富，而是为了回答问题。例如：

\begin{itemize}
  \item 误差是否集中在高峰期？
  \item 模型是否总是低估极端值？
  \item 某些区域是否因为训练数据稀少而表现更差？
  \item 数据清洗是否删除了真实但少见的样本？
  \item 高置信度预测是否真的更准确？
\end{itemize}

因此，绘图前应先写下要检查的问题。发现模式后，还要回到数据和代码中验证原因，不能仅凭图形作出结论。

\subsection{让AI辅助生成图形，但由人解释}

AI 很适合根据数据结构生成绘图代码，也可以建议应当查看哪些分组。但图形的含义仍需要结合领域背景解释。模型在某一区域误差较高，可能来自道路结构差异、传感器质量问题，也可能只是样本太少。AI 可以提出假设，最终判断仍需证据。

% ----------------------------------------------------------------
\section{组织一套完整验证流程}
% ----------------------------------------------------------------

五种方法分别解决不同问题。实际项目中，应将它们组织成有顺序的流程。

\subsection{第一步：确认评价对象}

明确当前要验证的是代码、数据处理结果、模型预测，还是根据模型生成的结论。不同对象需要不同证据。代码能够运行，不等于模型有效；模型分数较高，也不等于报告中的解释正确。

\subsection{第二步：建立基准和合理范围}

在查看复杂模型结果之前，确定简单基准、关键指标、合理数量级和不可违反的规则。这一步定义了“什么算值得采用”和“什么结果明显不合理”。

\subsection{第三步：运行常规评价并可视化错误}

在符合真实使用场景的数据划分上评价模型，报告总体指标、关键分组指标和错误图形。不要只保存最好分数，还应保存主要错误案例。

\subsection{第四步：测试关键极端情景}

根据任务风险选择少量重要情景，提前写出预期行为，再观察模型是否失效、如何失效，以及系统能否识别失效。

\subsection{第五步：反向追问并核对证据}

让 AI 或其他审查者主动寻找可能推翻结论的原因，并把每个怀疑转化为可执行检查。对于关键数字、引用和结论，应回到原始数据、实际运行结果或原始来源核对。

\subsection{第六步：形成带条件的评价结论}

最终结论不应只有“模型好”或“模型不好”，而应说明：

\begin{itemize}
  \item 模型在哪些场景优于基准；
  \item 哪些场景证据不足；
  \item 已经发现哪些主要错误模式；
  \item 使用时需要满足哪些条件；
  \item 出现什么情况时应停止使用。
\end{itemize}

下面是一段可以用于组织验证工作的 Prompt：

\begin{quote}
请为这个机器学习结果设计验证计划。依次给出：合理基准、单位与数量级检查、关键分组、三个重要极端情景、推荐的验证图形，以及可能推翻当前结论的证据。每项检查都要说明具体操作和判断标准。
\end{quote}

% ----------------------------------------------------------------
\section{案例：验证短时交通流预测结果}
% ----------------------------------------------------------------

假设 AI 建立了一个预测未来 15 分钟交通流量的模型，并报告其 RMSE 比历史均值低 25\%。

首先，评价者加入“上一时刻流量”和“上周同时刻流量”两个基准。结果显示，复杂模型比历史均值明显更好，但只比上一时刻基准改善 6\%。这说明原报告选择了较弱的基准，25\% 的提升夸大了模型的新增价值。

接着进行数量级检查。评价者发现少数路段的预测流量超过其历史最大值数倍。进一步检查后发现，这些路段的数据在表连接时被重复匹配。

修复数据后，评价者绘制真实值与预测值曲线，并分别计算高峰、平峰和暴雨时段的误差。模型在平峰表现较好，却经常低估高峰流量，在暴雨时段也没有明显优于简单基准。

最后，评价者模拟部分传感器缺失。模型在输入缺失时仍输出正常数值，却没有发出任何警告。于是最终结论被改写为：该模型在常规天气、数据完整的平峰时段可以提供参考；高峰、暴雨和传感器缺失时，应切换到基准方案或交由人工判断。

这个案例说明，验证并不是寻找一个漂亮分数，而是逐步确定结果在什么条件下正确、在什么条件下不应被采用。

% ----------------------------------------------------------------
\section{实战练习}
% ----------------------------------------------------------------

\subsection*{练习一：建立基准}

选择一个预测或分类任务，为它设计至少两个简单基准。说明每个基准代表什么已有规律，以及复杂模型需要达到什么条件才值得采用。

\subsection*{练习二：设计合理性检查表}

为一个 AI 生成的数据处理结果列出至少六项检查，包括单位、数量级、边界、组成关系和时间关系。为每项检查写出异常后的处理动作。

\subsection*{练习三：设计极端情景}

选择一个真实机器学习应用，设计三个极端情景。对每个情景说明输入如何变化、预期行为是什么，以及出现什么结果时应停止使用模型。

\subsection*{练习四：完整验证演练}

选择一个 AI 生成的模型结果，依次完成基准比较、合理性检查、关键分组评价、错误可视化、极端情景测试和反向追问。最后写出一段带有适用条件和停止条件的评价结论。
